\section{Билет 17. Волновые процессы. Типы волн. Продольные и поперечные волны. Волновая поверхность. Волновой фронт.}

\begin{definition}
\textbf{Волновым процессом} (или волной) называется процесс распространения колебаний в сплошной среде (твёрдой, жидкой или газообразной). При распространении волны частицы среды не перемещаются вместе с ней, а лишь колеблются около своих положений равновесия. Вместе с волной от частицы к частице передаётся состояние колебательного движения и его энергия. Основным свойством всех волн является \textbf{перенос энергии без переноса вещества}.
\end{definition}

\begin{definition}
В зависимости от физической природы возмущений выделяют основные \textbf{типы волн}:
\begin{enumerate}
    \item \textbf{Упругие (механические) волны} --- механические возмущения, распространяющиеся в упругой среде.
    \item Волны на поверхности жидкости.
    \item Электромагнитные волны.
\end{enumerate}
\end{definition}

\begin{definition}
Упругие волны классифицируются по направлению колебаний частиц относительно направления распространения волны:
\begin{itemize}
    \item \textbf{Продольные волны} --- частицы среды колеблются \textbf{вдоль} направления распространения волны. Возникают в средах, где возможны деформации сжатия и растяжения (твёрдые тела, жидкости, газы).
    \item \textbf{Поперечные волны} --- частицы среды колеблются в плоскостях, \textbf{перпендикулярных} направлению распространения волны. Возникают только в твёрдых телах, где возможны деформации сдвига.
\end{itemize}
\end{definition}

\begin{definition}
\textbf{Волновым фронтом} называется геометрическое место точек, до которых доходят колебания к данному моменту времени $t$.
\textbf{Волновой поверхностью} называется геометрическое место точек, колеблющихся в одинаковой фазе.
\end{definition}

\begin{remark}
Волновых поверхностей можно провести бесчисленное множество, а волновой фронт в каждый момент времени --- только один. При этом сам фронт также является волновой поверхностью. В зависимости от формы волны различают \textbf{плоские} (совокупность параллельных плоскостей), \textbf{сферические} (концентрические сферы) и \textbf{цилиндрические} волны. Форма волны определяется источником и свойствами среды.
\end{remark}
